% Generated draft: requires mathematical and source review.
\documentclass[11pt]{article}
\usepackage[margin=25mm]{geometry}
\usepackage{fontspec}
\setmainfont{Latin Modern Roman}
\usepackage{amsmath,amssymb,mathtools,unicode-math}
\setmathfont{Latin Modern Math}
\usepackage{enumitem,needspace,xurl,hyperref}
\setlength{\parindent}{0pt}
\setlength{\parskip}{5pt}
\setlength{\emergencystretch}{3em}
\newcommand{\N}{\mathbb{N}}
\newcommand{\Z}{\mathbb{Z}}
\newcommand{\Q}{\mathbb{Q}}
\newcommand{\R}{\mathbb{R}}
\newcommand{\C}{\mathbb{C}}
\newcommand{\F}{\mathbb{F}}
\newcommand{\A}{\mathbb{A}}
\newcommand{\PP}{\mathbb{P}}
\newcommand{\E}{\mathbb{E}}
\newcommand{\eps}{\varepsilon}
\newcommand{\dd}{\,\mathrm{d}}
\begin{document}
\section*{UnsolvedMath ID: 800616}
\textbf{Source problem code:} catalogue entry
\subsection{Definitions and mathematical statement}
Let \(\N=\{0,1,2,\ldots\}\), and give \(\Z^{\N}\) coordinatewise addition and the product topology of the discrete topology on \(\Z\). For an explicit conventional metric, put
\[
d(x,y)=\begin{cases}0,&x=y,\\2^{-\min\{n\in\N:x(n)\ne y(n)\}},&x\ne y.\end{cases}
\]
A subgroup \(G\le\Z^{\N}\) carries the induced topology and metric. It has \emph{strong measure zero} if, for every sequence \((\varepsilon_n)_{n\in\N}\) of positive real numbers, there are sets \(A_n\subseteq G\) such that
\[
G=\bigcup_{n\in\N}A_n\quad\text{and}\quad \operatorname{diam}_d(A_n)<\varepsilon_n\quad(n\in\N),
\]
where the diameter of the empty set is taken to be zero.

Write \(\mathcal O(G)\) for the family of open covers of \(G\). The property \(\mathsf S_{\mathrm{fin}}(\mathcal O,\mathcal O)\) means that for every sequence \((\mathcal U_n)_{n\in\N}\) in \(\mathcal O(G)\), there are finite subfamilies \(\mathcal F_n\subseteq\mathcal U_n\) with
\[
G=\bigcup_{n\in\N}\bigcup\mathcal F_n.
\]
The property \(\mathsf S_1(\mathcal O,\mathcal O)\) means that for every such sequence there are choices \(U_n\in\mathcal U_n\) satisfying \(G=\bigcup_{n\in\N}U_n\).

For an explicit eventual-domination interpretation, call \(G\) \emph{unbounded} if
\[
\forall h\in\N^{\N}\ \exists g\in G\quad
\bigl|\{n\in\N:|g(n)|>h(n)\}\bigr|=\infty.
\]
For subgroups this is equivalent to having no eventual upper bound in \(\Z^{\N}\) under the coordinatewise eventual order. This notion is not metric unboundedness.

With these conventions, the question is whether it is consistent with \(\mathsf{ZFC}\) that there exists a subgroup \(G\le\Z^{\N}\) which has strong measure zero, satisfies \(\mathsf S_{\mathrm{fin}}(\mathcal O,\mathcal O)\), is unbounded in the sense above, and fails \(\mathsf S_1(\mathcal O,\mathcal O)\). Explicitly, the last failure requires a sequence \((\mathcal U_n)_{n\in\N}\) of open covers of \(G\) for which every choice \(U_n\in\mathcal U_n\) satisfies \(\bigcup_n U_n\ne G\).

The supplied excerpt does not specify the metric, the meaning of unboundedness, or the background theory for consistency. The formulation above makes conventional choices explicit; their agreement with the intended source formulation requires review.
\subsection{Short English statement}
Can a model of set theory contain a subgroup of \(\Z^{\N}\) that has strong measure zero and the finite open-cover selection property, but is unbounded under eventual domination and fails the one-element open-cover selection property? The metric, unboundedness, and background-theory conventions used here require confirmation from the source context.
\subsection{Sources}
\begin{itemize}
\item Local UnsolvedMath registry, permanent ID \(800616\), exact source problem code \texttt{catalogue entry}; supplied title: \emph{Selection Principles in Mathematics --- Problem 44 ([BG])}. Upstream revision: \texttt{d371ee917ac9647fbb06cd6fc529deec6f3dda82}.
\item Tsaban, \emph{Selection Principles and Special Sets of Reals: Open Problems} (2006), Problem 44. Supplied source URL: \url{https://arxiv.org/abs/math/0609601}. The record identifies the extraction as \texttt{arxiv\_latex} and supplies the access date 2026-10-16; these are supplied provenance details, not independently verified access.
\item The registry lists Tsaban as proposer and 2006 as the proposed year, while the problem title and background retain the attribution/context marker \texttt{[BG]}. No expansion or bibliography entry for that marker is supplied.
\item The registry labels the problem \emph{open}. This is a supplied status label, not an independently verified literature assessment.
\end{itemize}
\end{document}
